Como podemos contar o número de ``furos'' em uma superfície? Para descrever e determinar este número formalmente, em sua tese de doutorado em 1851, Bernhard Riemann (1826-1866) definiu um número inteiro associado a certas supefícies $M$, \linebreak o \emph{gênero} $g$ de $M$, em termos de propriedades topológicas de $M$. Posteriormente, em uma série de artigos pioneiros em topologia algébrica, Henri Poincaré (1854-1912) desenvolveu um método para determinar $g$ em termos de contrações de curvas em $M$.%
\linebreak\vspace{-1.2em}

{\color{red} TODO motivação: $\pi_1(B^2) \iso 0$, $\pi_1(S^1) \iso \Z$, $\pi_1(T^2) \iso \pi(S^1) \times \pi(S^1) \iso \Z^2$}

\begin{comment}
Por exemplo, seja $D \subseteq \R^2$ um disco unitário e $S \subseteq \R^2$ um círculo unitário, ambos centrados na origem. (Veja a figura abaixo.) Dada uma curva fechada $\alpha : [0, 1] \to D$ passando por $p \in D$, é possível deformar $\alpha$ continuamente até obtermos uma curva $\alpha' : [0, 1] \to D$ constante no ponto $p$, i.e.\@ tal que $\alpha'(t) = p$ para todo $t \in [0, 1]$. Por outro lado, dada uma curva fechada $\alpha : [0, 1] \to S$ passando por $p \in S$, isto nem sempre é possível: em particular, pondo $p = (1,0)$ e $\alpha(t) = (\cos(t), \sen(t))$, isto é impossível. Intuitivamente, a primeira observação indica que $D$ ``não tem furos'', e portanto que $g = 0$, pois toda curva fechada pode ser contraída a um único ponto; enquanto a segunda indica que $S$ ``tem furos'', e portanto que $g > 0$, pois existem tais curvas que não podem ser continuamente deformadas em uma curva constante.%
\linebreak\vspace{-1.2em}

{\color{red} Figura: $D$ e $S$. Para fins ilustrativos, $S$ é um anel com espessura não-nula.}

{\color{red} Texto: relação entre $g$ e $\Omega(X, x_0) / \sim$}

Veremos nestas notas as definições e primeiros resultados em topologia algébrica acerca de \emph{caminhos}, \emph{homotopias} e o \emph{grupo fundamental} de um espaço, assim como noções preliminares e correlacionadas em teoria dos conjuntos, teoria dos grupos e topologia geral.  
\end{comment}